\chapter{The Needle and the Seam Rules}
\label{chap:needle-and-seam-rules}

The poison was a law of what the elaborator refuses. This chapter reads two more laws off the build, and
both are laws of economy --- of powers the construction cannot take for free, and of what it does instead.
The first is the rule of the needle: there is exactly one place the construction is allowed to identify two
things its own order cannot tell apart, and everything that must be so identified is routed through that one
place. The second is the rule of the seam, and it comes with a discovery: at the join between passes the
construction does not evaluate its reading but quotes it --- and the machine's own type discipline turned
out to enforce that quoting before the construction ever asked it to. Two laws, one theme: the construction
spends its costly powers once, at a single sanctioned place, or not at all.

\section{The needle: one truth told once}
\label{se:the-needle-one-truth-told-once}

There is one power the construction cannot take for free, and it is the power to identify. Telling two
things apart it can do everywhere, from its own decidable resources; but declaring that two things its order
cannot separate are nonetheless the same thing --- \emph{funging} the indistinguishable into one --- is a
step outside the order, and a step outside the order cannot be taken casually or the whole discipline
dissolves. So the construction takes it exactly once. There is a single sanctioned identification, one act
of quotient-soundness, and every place in the whole work that must treat two indistinguishable readings as
one is routed through that single needle. Nothing gets a second. The rule extends to the smallest scale: a
value the construction needs twice is bound once and used twice, never computed twice, so that one truth is
told exactly once and every later use points back to that one telling. Write an expression a second time and
you have made a second thing the elaborator must check is the same as the first --- a second needle, an
identification not sanctioned. The construction refuses to; it threads the one needle and passes everything
through it. This is why there is only ever one: the genuine identification of the indistinguishable is the
one power that cannot be had for nothing, and a power that costs is spent once and remembered, not scattered.

\section{The seam: quote, do not evaluate}
\label{se:the-seam-quote-not-evaluate}

At a seam --- the join where one pass hands the next its residue --- the construction faces a proposition it
might try to settle: whether the two readings the seam pairs agree. It does not settle it. Instead of
running the decision at the seam, it \emph{quotes} the proposition --- places it into the cell of the tape
unevaluated, as a statement held and not yet judged --- and carries it forward as content. The decision is
not made where the seam is written; it is deferred to the place the cell is later read, where the two facts
the proposition compares supply their own decidability and settle it by a plain case-split. The seam records
what is to be decided; the reading decides it. Quoting rather than evaluating is not a stylistic preference
here but the construction's method at every join: a seam is a place to write down a question, not to answer
it, and the answer waits for a reader who has what it takes to give one.

\section{The law the machine enforced first}
\label{se:law-machine-enforced-first}

Here is the discovery, and it is the reason this rule is a law and not a habit: the machine enforced
quote-don't-evaluate before the construction asked it to. The cell that holds the seam's proposition holds
it in a slot that is a bare proposition --- and a bare proposition, in the machine's own type discipline,
has no decision attached to it. It cannot be asked to evaluate, because there is nothing there to run: a
proposition is a statement of what would be true, not a procedure for finding out. So the construction, in
trying to place the seam's question into the tape, found that the only thing the slot would accept was the
question itself, unevaluated --- the decision simply could not be demanded there, because the type of the
slot had no decision to give. What might have been imposed as a discipline was instead discovered as a
constraint: the machine would not let the seam evaluate, and the construction did not need it to. The
decision belongs at the read, where the facts carry their own way of settling it; at the seam there is only
the quoted question, and the type of the slot is what guarantees it stays that way. A rule that would have
been a good idea turned out to be a fact of the material --- the strongest kind of law, one the construction
does not have to keep because the medium keeps it.

\section{What the two rules share}
\label{se:what-the-two-rules-share}

The needle and the seam are one economy seen twice. Both concern a power the construction must not spend
loosely --- the needle, the power to identify what cannot be separated; the seam, the power to decide a
proposition --- and both resolve the same way: the costly act is taken at exactly one sanctioned place, or
deferred to where it comes for free. The identification is spent once, at the single needle, and never
again. The decision is not spent at the seam at all but postponed to the read, where the facts pay for it
themselves with their own decidability, and the seam \emph{tanges} nothing --- it only records. What is
established is that the construction's discipline about its two irreducible powers is not imposed from
taste but enforced by the elaborator: the identification is threaded once because a second would be an
unsanctioned quotient, and the decision is quoted at the seam because the slot's own type will hold nothing
else. The economy is the machine's, read off what it would and would not permit. The next law is of the
tape itself --- that it only ever grows --- and of the certificates the construction reads by exhaustion.
